\chapter{1977年陕西卷}

\begin{problem}
求值：\( \left( \frac{1}{300} \right)^{-\frac{1}{2}}+10\left( \frac{\sqrt{3}}{2} \right)^{\frac{1}{2}}\cdot\left( \frac{27}{4} \right)^{\frac{1}{4}}-10\left( 2-\sqrt{3} \right)^{-1} \).
\end{problem}

\begin{solution}
先化简各项，有
\[
\left( \frac{1}{300} \right)^{-\frac{1}{2}}=\sqrt{300}=10\sqrt{3}
\]
又
\[
\left( \frac{\sqrt{3}}{2} \right)^{\frac{1}{2}}\left( \frac{27}{4} \right)^{\frac{1}{4}}=\frac{\sqrt[4]{3}}{\sqrt{2}}\cdot\frac{\sqrt[4]{3^3}}{\sqrt{2}}=\frac{3}{2}
\]
由于 \(\left(2-\sqrt{3}\right)\left(2+\sqrt{3}\right)=1\)，所以 \(\left(2-\sqrt{3}\right)^{-1}=2+\sqrt{3}\). 因此原式为
\[
10\sqrt{3}+10\cdot\frac{3}{2}-10\left(2+\sqrt{3}\right)=10\sqrt{3}+15-20-10\sqrt{3}=-5
\]
\end{solution}

\begin{problem}
求值：\( \sin 480^\circ\cdot\cos 480^\circ\cdot\cos 960^\circ\cdot\cos 1920^\circ \).
\end{problem}

\begin{solution}
利用三角函数的周期性，\(480^\circ=360^\circ+120^\circ\)，\(960^\circ=2\cdot360^\circ+240^\circ\)，\(1920^\circ=5\cdot360^\circ+120^\circ\). 因而
\[
\begin{gathered}
\sin480^\circ=\sin120^\circ=\frac{\sqrt{3}}{2},\qquad \cos480^\circ=\cos120^\circ=-\frac{1}{2},\\
\cos960^\circ=\cos240^\circ=-\frac{1}{2},\qquad \cos1920^\circ=\cos120^\circ=-\frac{1}{2}.
\end{gathered}
\]
所以原式为
\[
\frac{\sqrt{3}}{2}\left(-\frac{1}{2}\right)^3=-\frac{\sqrt{3}}{16}
\]
\end{solution}

\begin{problem}
化简：\( \left| 6-a \right|-\left| 2a+1 \right|+\sqrt{a^2+10a+25} \)（\( a<-5 \)）.
\end{problem}

\begin{solution}
由 \(a<-5\)，可得 \(6-a>0\)、\(2a+1<0\)、\(a+5<0\). 又
\(a^2+10a+25=(a+5)^2\)，所以
\[
\left|6-a\right|-\left|2a+1\right|+\sqrt{a^2+10a+25}=(6-a)+(2a+1)-\left(a+5\right)=2
\]
\end{solution}

\begin{problem}
解不等式：\( \lg \left( 2x+1 \right)>\lg \left( 5-x \right) \).
\end{problem}

\begin{solution}
对数的真数必须为正，因此先得
\(2x+1>0\)，\(5-x>0\)
即 \(x>-\frac{1}{2}\) 且 \(x<5\). 在此定义域内，常用对数函数单调递增，故原不等式等价于
\[
2x+1>5-x,\qquad\text{即}\qquad x>\frac{4}{3}
\]
与定义域取交集，得到原不等式的解集为
\[
\frac{4}{3}<x<5
\]
\end{solution}

\begin{problem}
设 \(A\)，\(B\) 为锐角，且 \( \sin A=\frac{3}{5} \)，\( \sin B=\frac{5}{13} \)，求 \( \tan\left( A-B \right) \).
\end{problem}

\begin{solution}
因为 \(A\)、\(B\) 均为锐角，故余弦值为正，从而
\(\cos A=\sqrt{1-\sin^2A}=\sqrt{1-\left(\frac{3}{5}\right)^2}=\frac{4}{5}\)，\(\cos B=\sqrt{1-\sin^2B}=\sqrt{1-\left(\frac{5}{13}\right)^2}=\frac{12}{13}\).
由两角差公式，
\[
\sin(A-B)=\frac{3}{5}\cdot\frac{12}{13}-\frac{4}{5}\cdot\frac{5}{13}=\frac{16}{65}
\]
\[
\cos(A-B)=\frac{4}{5}\cdot\frac{12}{13}+\frac{3}{5}\cdot\frac{5}{13}=\frac{63}{65}>0
\]
因此
\[
\tan(A-B)=\frac{\sin(A-B)}{\cos(A-B)}=\frac{16/65}{63/65}=\frac{16}{63}
\]
\end{solution}

\begin{problem}
实数 \( p \)，\( q \) 应满足怎样的条件，才能使方程 \( x^2+px+q=0 \) 的两根成为一直角三角形两锐角的正弦？
\end{problem}

\begin{solution}
设直角三角形的一个锐角为 \(\alpha\)，则 \(0<\alpha<\frac{\pi}{2}\)，另一个锐角为 \(\frac{\pi}{2}-\alpha\). 若方程的两根为 \(x_1\)、\(x_2\)，则可取
\(x_1=\sin\alpha\)，\(x_2=\sin\left(\frac{\pi}{2}-\alpha\right)=\cos\alpha\).
由韦达定理，
\(\sin\alpha+\cos\alpha=-p\)，\(\sin\alpha\cos\alpha=q\).
将第一个等式平方，并利用 \(\sin^2\alpha+\cos^2\alpha=1\)，得到
\[
p^2=\left(\sin\alpha+\cos\alpha\right)^2=1+2\sin\alpha\cos\alpha=1+2q
\]
此外，\(0<\alpha<\frac{\pi}{2}\) 给出
\(1<\sin\alpha+\cos\alpha\leq\sqrt{2}\)，\(0<\sin\alpha\cos\alpha\leq\frac{1}{2}\).
因此必要条件为
\(1+2q=p^2\)，\(-\sqrt{2}\leq p<-1\)，\(0<q\leq\frac{1}{2}\).

再证这些条件充分. 设 \(s=-p\)，则 \(1<s\leq\sqrt{2}\)，且由关系式 \(q=\frac{s^2-1}{2}\). 原方程可写为
\[
x^2-sx+q=0
\]
它的判别式为
\[
\Delta=s^2-4q=s^2-2\left(s^2-1\right)=2-s^2\geq0
\]
设两根为 \(u\)、\(v\)，则 \(u+v=s>0\)、\(uv=q>0\)，所以 \(u\)、\(v\) 均为正数；又
\[
u^2+v^2=(u+v)^2-2uv=s^2-2q=1
\]
故 \(0<u<1\)、\(0<v<1\). 取锐角 \(\alpha\) 使 \(\sin\alpha=u\)，则 \(\cos\alpha=\sqrt{1-u^2}=v\)，所以两根确为 \(\sin\alpha\)、\(\cos\alpha\)，即为直角三角形两个锐角的正弦.

所以所求条件是
\(1+2q=p^2\)，\(-\sqrt{2}\leq p<-1\)，\(0<q\leq\frac{1}{2}\).
\end{solution}

\begin{problem}
以任意一点到已知三角形三个顶点连线的中点为顶点作三角形. 所作三角形的面积是已知三角形的几倍？
\end{problem}

\begin{solution}
如图，\(P\) 为已知 \(\triangle ABC\) 外任一点，\(A'\)、\(B'\)、\(C'\) 分别是 \(PA\)、\(PB\)、\(PC\) 的中点.
\solutionfigure{\bitmapfigure[max width=5.2cm]{img/1977/shaanxi/q7_fig1.png}}
由中点的向量表示，
\(\overrightarrow{A'B'}=\frac{1}{2}\overrightarrow{AB}\)，\(\overrightarrow{B'C'}=\frac{1}{2}\overrightarrow{BC}\)，\(\overrightarrow{A'C'}=\frac{1}{2}\overrightarrow{AC}\).
因此
\[
\frac{A'B'}{AB}=\frac{B'C'}{BC}=\frac{A'C'}{AC}=\frac{1}{2}
\]
从而 \(\triangle A'B'C'\sim\triangle ABC\). 相似三角形的面积比等于相似比的平方，所以
\[
\frac{S_{\triangle A'B'C'}}{S_{\triangle ABC}}=\left(\frac{1}{2}\right)^2=\frac{1}{4}
\]
故所作三角形的面积是已知三角形面积的 \(\frac{1}{4}\). 这个结论与 \(P\) 在三角形外、内或边上的位置无关.
\end{solution}

\begin{problem}
求过椭圆 \( \frac{x^2}{13^2}+\frac{y^2}{12^2}=1 \) 的右侧的焦点，且与圆 \( x^2+y^2=3^2 \) 相切的直线方程.
\end{problem}

\begin{solution}
椭圆的半长轴、半短轴分别为 \(a=13\)、\(b=12\)，故焦半距为
\[
c=\sqrt{a^2-b^2}=\sqrt{13^2-12^2}=5
\]
所以右焦点为 \(F(5,0)\).

设所求切线与圆的切点为 \(T(x_1,y_1)\). 圆 \(x^2+y^2=9\) 在 \(T\) 处的切线方程为
\[
x_1x+y_1y=9
\]
由于切线经过 \(F(5,0)\)，有 \(5x_1=9\)；又因为 \(T\) 在圆上，有 \(x_1^2+y_1^2=9\). 解得
\(x_1=\frac{9}{5}\)，\(y_1^2=9-\left(\frac{9}{5}\right)^2=\frac{144}{25}\)，\(y_1=\pm\frac{12}{5}\).
分别代入切线方程，得到
\(9x+12y-45=0\)，\(9x-12y-45=0\).
将两式分别约去公因数 \(3\)，所求直线为
\(3x+4y-15=0\)，\(3x-4y-15=0\).
\end{solution}

\begin{problem}
三角形的一内角为 \( 30^\circ \)，它的一邻边长为 \( 4 \)，对边长为 \( \frac{5}{2} \)，求另一边长.
\end{problem}

\begin{solution}
设所求另一边长为 \(x\). 取已知的 \(30^\circ\) 为两已知邻边 \(4\)、\(x\) 的夹角，则由余弦定理，
\[
\left(\frac{5}{2}\right)^2=x^2+4^2-2\cdot x\cdot4\cos30^\circ=x^2+16-4\sqrt{3}x
\]
两边同乘以 \(4\)，整理得
\[
4x^2-16\sqrt{3}x+39=0
\]
其判别式为
\[
\Delta=\left(-16\sqrt{3}\right)^2-4\cdot4\cdot39=144=12^2
\]
所以
\[
x=\frac{16\sqrt{3}\pm12}{8}=\frac{4\sqrt{3}\pm3}{2}
\]
两值均为正. 将任一正根作为一条邻边，与长度为 \(4\) 的邻边夹角取 \(30^\circ\)，即可构成非退化三角形；由上面的余弦定理，该三角形的对边正好为 \(\frac{5}{2}\)，所以两根均不需舍去. 故另一边长为
\[
\frac{4\sqrt{3}+3}{2}\quad\text{或}\quad\frac{4\sqrt{3}-3}{2}
\]
\end{solution}

\begin{problem}
在一块边长为 \( a \) 的正方形上，切着两边截去了一个半径为 \( r \) 的圆料（\( r\le\frac{\left( 2-\sqrt{2} \right)a}{4} \)）.

\begin{enumerate}
\item 如果象图 1 那样截出一个正方形的料，它的面积为多大？
\item 如果象图 2 那样截出的一个正方形的料，它的面积为多大？这二种截法截下的正方形的料的面积哪个较大？
\end{enumerate}

\bitmapfigure[max width=4.6cm]{img/1977/shaanxi/q10_fig1.png}\hfill
\bitmapfigure[max width=4.6cm]{img/1977/shaanxi/q10_fig2.png}
\end{problem}

\begin{solution}
\begin{enumerate}
\item 如图 1. 设圆 \(\odot O\) 与正方形 \(ABCD\) 的边 \(AB\)、\(BC\) 分别相切于 \(E\)、\(F\). 则 \(OEBF\) 是正方形，且圆心 \(O\) 在对角线 \(BD\) 上. 因此
\(BD=\sqrt{2}a\)，\(OB=\sqrt{2}r\).
由于 \(B'\) 是对角线 \(BD\) 与圆的交点，且 \(O\) 位于 \(B\)、\(B'\) 之间，\(OB'=r\)，所以
\(BB'=BO+OB'=\sqrt{2}r+r\)，\(qquad B'D=BD-BB'=\sqrt{2}a-(\sqrt{2}+1)r\).
正方形 \(A'B'C'D\) 的对角线为 \(B'D\)，故其面积为
\[
S_1=\frac{1}{2}(B'D)^2=\frac{1}{2}\left[\sqrt{2}a-(\sqrt{2}+1)r\right]^2=\left[a-\left(1+\frac{\sqrt{2}}{2}\right)r\right]^2
\]

\item 如图 2. 设圆 \(\odot O\) 与 \(AB\)、\(BC\) 分别相切于 \(E\)、\(F\)，并与正方形 \(A'B'C'D'\) 的边 \(A'B'\) 相切于 \(G\). 由于 \(A'B'=A'D'\)，\(A'B'\perp A'D'\)，且 \(AB\perp AD\)，有 \(\angle AD'A'=\angle BA'B'\). 两个相应直角三角形全等，从而
\[
AA'=BB'=r+FB'
\]
从同一点 \(A'\) 向圆作的两条切线段相等，故 \(A'G=A'E\). 于是
\[
A'G=A'E=AB-AA'-EB=a-(r+FB')-r=a-2r-FB'
\]
同理，从点 \(B'\) 作圆的两条切线段相等，故 \(B'G=B'F\). 因而
\[
A'B'=A'G+B'G=a-2r-FB'+FB'=a-2r
\]
所截正方形的面积为
\[
S_2=(A'B')^2=(a-2r)^2
\]

因为 \(r>0\) 且 \(1+\frac{\sqrt{2}}{2}<2\)，又由题设可知两边长均为正，所以
\[
a-\left(1+\frac{\sqrt{2}}{2}\right)r>a-2r>0
\]
平方后得到 \(S_1>S_2\). 因此第一种截法截下的正方形料面积较大.
\end{enumerate}
\end{solution}

\section{参考题}

\begin{problem}
在一个椭圆的所有内接矩形中，怎样的矩形面积最大？

\bitmapfigure[max width=6.2cm]{img/1977/shaanxi/q11_fig1.png}
\end{problem}

\begin{solution}
设椭圆方程为 \(\frac{x^2}{a^2}+\frac{y^2}{b^2}=1\)，其中 \(a>b>0\). 先说明任意内接矩形的中心都是椭圆中心. 将坐标按 \(X=\frac{x}{a}\)、\(Y=\frac{y}{b}\) 作伸缩变换，椭圆变为单位圆，原矩形变为一个内接平行四边形. 圆内接平行四边形必为矩形，其对角线为直径，所以该平行四边形的中心是圆心，从而原矩形的中心是椭圆中心.

设该矩形中心为原点，取从中心指向相邻顶点的两个半边向量的坐标分别为 \((u_1,u_2)\)、\((v_1,v_2)\). 因为矩形的两边垂直，
\[
u_1v_1+u_2v_2=0
\]
矩形的四个顶点为两向量的带符号和，且都在椭圆上. 比较点 \((u_1+v_1,u_2+v_2)\) 与 \((u_1-v_1,u_2-v_2)\) 的椭圆方程，得到
\[
\frac{u_1v_1}{a^2}+\frac{u_2v_2}{b^2}=0
\]
又因 \((v_1,v_2)=\lambda(-u_2,u_1)\)，其中 \(\lambda\ne0\)，代入上式可得
\[
\lambda u_1u_2\left(\frac{1}{b^2}-\frac{1}{a^2}\right)=0
\]
由于 \(a>b>0\)，所以 \(u_1u_2=0\). 因而矩形的两边分别平行于椭圆的两条坐标轴.

设矩形的半长为 \(x\)，半宽为 \(y\)，则顶点 \((x,y)\) 在椭圆上，故
\(\frac{x^2}{a^2}+\frac{y^2}{b^2}=1\)，\(qquad y=\frac{b}{a}\sqrt{a^2-x^2}\)，\(qquad 0\leq x\leq a\).
于是矩形面积为
\[
S(x)=4xy=\frac{4b}{a}x\sqrt{a^2-x^2}=\frac{4b}{a}\sqrt{x^2\left(a^2-x^2\right)}
\]
由算术几何平均不等式，
\[
\sqrt{x^2\left(a^2-x^2\right)}\leq\frac{x^2+a^2-x^2}{2}=\frac{a^2}{2}
\]
所以 \(S(x)\leq2ab\). 当且仅当 \(x^2=a^2-x^2\)，即 \(x=\frac{\sqrt{2}}{2}a\) 时取等号；此时
\(2x=\sqrt{2}a\)，\(2y=\sqrt{2}b\)，\(qquad S_{\max}=2ab\).
因此，当内接矩形的长、宽分别为椭圆半长轴、半短轴的 \(\sqrt{2}\) 倍时，面积最大.
\end{solution}

\begin{problem}
试解
\( \frac{\mathrm{d} y}{\mathrm{d} x}-2xy=2x^3 \)，\( y\big|_{x=0}=1 \).
\end{problem}

\begin{solution}
原方程是一阶线性微分方程. 取积分因子 \(\e^{-x^2}\)，因为
\[
\frac{\mathrm{d}}{\mathrm{d}x}\e^{-x^2}=-2x\e^{-x^2}
\]
所以方程两边同乘以该积分因子后，左边恰为一个乘积的导数：
\[
\frac{\mathrm{d}}{\mathrm{d}x}\left(y\e^{-x^2}\right)=2x^3\e^{-x^2}
\]
令 \(u=x^2\)，则 \(\mathrm{d}u=2x\,\mathrm{d}x\)，从而
\[
\int2x^3\e^{-x^2}\,\mathrm{d}x=\int u\e^{-u}\,\mathrm{d}u=-(u+1)\e^{-u}+C=-(x^2+1)\e^{-x^2}+C
\]
因此
\[
y\e^{-x^2}=-(x^2+1)\e^{-x^2}+C
\]
即原方程的通解为
\[
y=C\e^{x^2}-x^2-1
\]
代入初值条件 \(y(0)=1\)，得 \(C-1=1\)，所以 \(C=2\). 故初值问题的解为
\[
y=2\e^{x^2}-x^2-1
\]
直接求导可得 \(y'=4x\e^{x^2}-2x\)，代回原方程有 \(y'-2xy=2x^3\)，且 \(y(0)=1\)，故该解满足题设.
\end{solution}
